\documentclass[onecolumn, draftclsnofoot,10pt, compsoc]{IEEEtran}

\usepackage[doublespacing]{setspace}
\usepackage[margin=0.75in]{geometry}
\usepackage{graphicx}
\usepackage{hyperref}
\usepackage[round, sort, numbers]{natbib}

\hbadness=1000 % suppress warnings

\begin{document}
	\begin{titlepage}

	\center

	\textsc{\LARGE Oregon State University}\\[1.5cm]
	\textsc{\Large CS461 Fall 2017 }\\[0.5cm]
	\textsc{\large Group 19}\\[0.5cm] % Minor heading such as course title

	\vfill

	{\huge\bfseries Aerolyzer Library Technology Review}\\[0.4cm]

	\vfill

	\begin{minipage}{0.4\textwidth}
		\begin{flushleft}
			\large
			\textit{Authors}\\
			Kin-Ho \textsc{Lam}\\
			Logan \textsc{Wingard}\\
			Daniel \textsc{Ross}\\
		\end{flushleft}
	\end{minipage}
	~
	\begin{minipage}{0.4\textwidth}
		\begin{flushright}
			\large
			\textit{Supervisor}\\
			Dr. Kim \textsc{Whitehall}
		\end{flushright}
	\end{minipage}

	\vfill

	\begin{abstract}
		\begin{singlespace}
			This document shall discuss possible technologies the Aerolyzer project shall leverage to accomplish horizon detection, aerosol content inference, and data regression analysis.
		\end{singlespace}
	\end{abstract}

	\vfill
	\vfill
	\vfill
	{\large11/12/17}
	\vfill
\end{titlepage}

\section{Table of Contents}
\tableofcontents
\clearpage

\begin{singlespace}

\section{Sunset and Sunrise Classification}
	\subsection{Overview}
		Atmospheric aerosols are particles suspended in the atmosphere.
		These particles originate from natural sources such as volcanic eruptions and unnatural sources such as fossil fuel combustion.
		One can visibly see the effects of aerosols in the 'Rayleigh scattering effect' which visibly reddens sunsets and sunrises.
		The objective of the Aerolyzer project is to create a tool that infers air quality by performing analysis on the distribution of sunlight hues in images of sunsets and sunrises.
		Before aerosol analysis can be performed, the Aerolyzer library must first filter out images that cannot be analyzed.
		Unacceptable images include pictures that do not feature a sunset, sunrise, or have objects obscuring too much of the horizon. 
		However, variables such as camera manufacturer, perspective, post-processing, image format, obstructions, and lighting prevents a simple algorithm from classifying all possible variations.
	
	\subsection{Criteria}
		Aerolyzer's unique image criteria necessitates a predictive data model capable of detecting horizons with a minimum of 66\% certainty, and sunsets or sunrises with a minimum of 50\% certainty.
		Third-party tools or libraries to be used in Aerolyzer's horizon classifier shall be compared by availability and utility.
		The Aerolyzer project shall ideally use free and non-proprietary libraries to perform image analysis.
		Such libraries are preferably open source with usage licenses compatible with \href{https://www.apache.org/licenses/LICENSE-2.0}{Apache License 2.0}.

		\begin{center}
			\begin{tabular}{|l|p{3cm}|p{4cm}|p{5cm}|}
			\hline \textbf{} & \textbf{Cost} & \textbf{License} & \textbf{Tools/Advantages} \\\hline
			OpenCV & Free & Open Source, BSD License & Robust computer-vision library\\\hline
			Google Tensorflow & Free & Creative Commons Attribution 3.0 License, Apache License 2.0 & Supervised machine learning library with proven image-recognition modules\\\hline
			Amazon AWS Rekognition & \$0.01/1,000 images & Proprietary \& Closed source& Supervised machine learning library with cloud platform processing\\\hline
			\end{tabular}
		\end{center}

	\subsection{OpenCV}
		OpenCV is a BSD-licensed computer-vision library with many relevant sub-modules.
		Containing over 2,500 algorithms, OpenCV is capable of analyzing images and videos to recognize faces, identify objects, perform visual transformations, and much more.
		Open ALPR is an example of a powerful image classifier built with OpenCV.
		Open APLR is an automatic license plate recognition library, capable of identifying and reading license plates and car models in near real-time. \cite{matthill}
		The functions for distinguishing cars and license plates is conceptually similar to Aerolyzer's horizon detection goals.
		OpenCV's \textit{imgproc} module performs image filtering, geometrical image transformations, and color space conversion. 
		The \textit{objdetect} module is able to detect predefined classes such as faces, eyes, and cars.
		These modules present useful tools that are capable of horizon detection given proper tolerances and algorithms.
		One can use image transformations to normalize image data such that a classifier can more easily distinguish what is in the picture.
		Pre-defined classes are useful in detecting obstructions that may block a good view of the horizon.
		OpenCV is a potential candidate for horizon/sunset/sunrise detection due to its feature-rich library, ample documentation, and pedigree of successful image classifiers.

	\subsection{Tensorflow}
		Tensorflow is a supervised machine learning library created and distributed by Alphabet \textsuperscript{Inc.}'s subsidiary Google.
		The library is broad and contains tools to accurately characterize a wide variety of potential weather conditions provided an adequate data set.
		One of the leading machine learning libraries, Tensorflow has been proven to build versatile and high accuracy predictive data models.
		Rohan Verma's Galaxy Image Classifier is an example of a powerful classification library built with Tensorflow. \cite{rhnvrm}
		This library is capable of detecting whether a submitted image contains a spiral or elliptical galaxy.
		The process to create Verma's galaxy classifier is conceptually similar to building a predictive model for images of horizons, sunsets, or sunrises.
		Tensorflow is best leveraged in conjunction with OpenCV to take advantage of both supervised and unsupervised machine learning.
		

	\subsection{Amazon AWS Rekognition}
		Amazon Rekognition is a supervised machine learning library with a focus on usability.
		Rekognition performs analysis on Amazon’s servers, its API is set up to service image recognition calls on-demand.
		When provided an image, Rekognition will simply return a JSON data structure containing elements detected in the image as well as a confidence score.
		Rekognition is a less preferable alternative to OpenCV and Tensorflow due to its processing cost, limited library, and closed-source proprietary algorithm.		

	\subsection{Conclusion}
		Although the final algorithm is subject to change; at the request of Dr. Whitehall, horizon classification development shall focus on OpenCV algorithms.
		Tensorflow and OpenCV seem to be the strongest candidates to build a predictive model to filter unacceptable images.

\clearpage

\section{Aerosol Inference}
	\subsection{Overview}
		Monitoring atmospheric aerosols is important due to their effects on the atmosphere's chemical composition and radiation distribution.
		The Aerolyzer library hopes to feature a module that infers aerosol size based on the hue distribution in acceptable horizon images.
		Creating such a library will require incorporating machine learning and statical libraries to perform regression and hypothesis testing.
		
	\subsection{Criteria}
		Third-party tools or libraries to be used in Aerolyzer's aerosol inference module shall be compared by utility.
		Desirable libraries shall incorporate statistical analysis or machine learning to identify patterns and analyze trends.

		\begin{center}
			\begin{tabular}{|l|p{3cm}|p{5cm}|p{5cm}|}
			\hline \textbf{} & \textbf{Cost} & \textbf{License} & \textbf{Tools/Advantages} \\\hline
			Google Tensorflow & Free & Creative Commons Attribution 3.0 License, Apache License 2.0 & Supervised machine learning library with proven regression-analysis modules\\\hline
			TF Learn & Free & Creative Commons Attribution 3.0 License, Apache License 2.0 & High-level API wrapper for Tensorflow\\\hline
			OpenCV & Free & Open Source, BSD License & Support Vector Machine library\\\hline
			\end{tabular}
		\end{center}

	\subsection{Google Tensorflow}
		Google Tensorflow supports deep-learning data regression.
		This enables an individual to train Tensorflow's supervised machine learning algorithm to recognize and predict trends by following models such as drawing the line of best fit in a given set of data.
		One would apply Tensorflow to the Aerolyzer project by feeding the algorithm hundreds of images which features aerosols of known sizes.
		A large quantity of training images is necessary to improve accuracy.
		After given a large selection of example images, Tensorflow will generate a classifier.
		When the generated Tensorflow classifier is fed an image, it will give its approximation as a confidence score as to how well the image aligns with its training data.

	\subsection{TFLearn}
		TFLearn is a wrapper built around Google Tensorflow.
		It features a modular deep learning library with an easy-to-use higher-level API to enable faster algorithm prototyping.
		While the underlying algorithms are still Tensorflow, TFLearn is advantageous in that it simplifies Tensorflow functions into less verbose calls. 

	\subsection{OpenCV}
		OpenCV's Support Vector Machine (SVM) library is a supervised machine learning algorithm.
		Given a sample dataset, SVM generates an optimized hyperplane model to categorize future input.
		For example, an SVM can be trained by providing several pre-sorted sets of images of sunsets and sunrises.
		This can be applied to the Aerolyzer project by training an SVM with a set of horizons with known aerosol content.

	\subsection{Conclusion}
		While the overall software focus is subject to change throughout the development of the Aerolyzer library, some supervised machine learning library is an optimum choice for performing aerosol analysis.
		The wide variation of inputs makes designing a simple horizon hue classifier inadequate for approximating aerosol size.

\clearpage

\section{Data Interpretation}
	\subsection{Overview}
		Aerosol data aggregation is largely unreadable to individuals who do not have in-depth atmospheric knowledge.
		The Aerolyzer library requires a machine-learning tool to 

	\subsection{Criteria}
	The Aerolyzer project will need to generate a trend analysis using the data it extrapolates. 
			
		\begin{center}
			\begin{tabular}{|l|p{3cm}|p{5cm}|p{5cm}|}
			\hline \textbf{} & \textbf{Cost} & \textbf{License} & \textbf{Tools/Advantages} \\\hline
			R \& Matlab Libraries & Free & Apache License 2.0 & \\\hline
			Apache Hadoop & Free & Apache License 2.0 & Scales R \& Matlab to work with large sets of data\\\hline
			& & & \\\hline
			\end{tabular}
		\end{center}

	\subsection{R \& Matlab Libraries}
	Leveraging the statistical power of R and Matlab libraries will enable the Aerolyzer project to interpolate weather patterns from sampled images.
	This shall be accomplished through regression analysis which shall consider dates of when images are taken, the inferred aerosol content of the image, patterns referenced from 3rd party APIs.
	R and Matlab libraries are readily available modules in the Python library which can be applied to Aerolyzer’s future dataset. 

	\subsection{Conclusion}


	\clearpage
	\nocite{*}
	\bibliographystyle{plain}
	\bibliography{ref}

\end{singlespace}
\end{document}
